\chapter{Контрольная по занятиям №1--3}\label{mock13:start}
\textbf{90 минут, 100 баллов.} Вариант охватывает лекции и семинары №1--3. На решение отводится примерно 80 минут, на проверку --- 10 минут. Работайте на бумаге без калькулятора, программ, учебных материалов и ответов. Показывайте промежуточные вычисления. Баллы предназначены для самопроверки.

\textbf{Соглашения.} В заданиях В2--В5 ширина целочисленных слов равна 8 битам, кроме явно указанных случаев. Знаковые значения представлены дополнительным кодом, если не назван другой формат. Биты нумеруются справа с нуля. Флаги записываются в порядке CF, OF, ZF, SF; для вычитания используются правила SUB x86. Подпункты независимы.

В заданиях В6--В8 используется IEEE 754 \texttt{float32}: 1 бит знака, 8 битов порядка, 23 бита дробной части, bias $=127$. Округление --- к ближайшему представимому числу; при равных расстояниях выбирается число с чётной значащей частью, то есть с нулём в её младшем сохраняемом бите. Каждая арифметическая операция округляется отдельно до \texttt{float32}. Hex-запись задаёт 32-битное слово целиком, старшими разрядами слева.

\section{В1. Переводы и арифметика (10 баллов, 8 минут)}\label{task:v1}
\begin{enumerate}
    \item Переведите $205_{10}$ в двоичную и шестнадцатеричную системы. Покажите способ перевода. \textbf{(6 баллов)}
    \item Вычислите $57_8+26_8$ с указанием переносов. Проверьте ответ в десятичной системе. Разрядность не ограничена. \textbf{(4 балла)}
\end{enumerate}

\section{В2. Диапазоны и коды (10 баллов, 10 минут)}\label{task:v2}
\begin{enumerate}
    \item Укажите диапазон и количество различных значений \textbf{7-битного} дополнительного кода. \textbf{(2 балла)}
    \item Представьте $-27$ в \textbf{8-битных} прямом, обратном, дополнительном кодах и коде со смещением 128. Для каждого кода покажите получение, двоичную и hex-записи. \textbf{(8 баллов)}
\end{enumerate}

\section{В3. Арифметика и флаги (15 баллов, 10 минут)}\label{task:v3}
Выполните операции в 8 битах. Для каждой укажите hex-код результата, его знаковое значение и четыре флага. Обоснуйте CF и OF.
\begin{enumerate}
    \item $85+60$. \textbf{(7 баллов)}
    \item $40-75$. Получите биты результата через сложение с дополнительным кодом вычитаемого. Флаги укажите для исходного вычитания. \textbf{(8 баллов)}
\end{enumerate}

\section{В4. Побитовые операции (10 баллов, 8 минут)}\label{task:v4}
Для исходного 8-битного $x=B3_{16}$ независимо выполните:
\begin{enumerate}
    \item $x$ AND $5C_{16}$. \textbf{(2 балла)}
    \item Сброс бита 5. \textbf{(2 балла)}
    \item Установку бита 2. \textbf{(2 балла)}
    \item Извлечение битов с номерами 2--4 как unsigned-числа от 0 до 7. \textbf{(4 балла)}
\end{enumerate}
В первом подпункте покажите побитовое вычисление. В остальных укажите операции, маски и результаты; в четвёртом --- также сдвиг и десятичное значение поля.

\section{В5. Сдвиги (15 баллов, 8 минут)}\label{task:v5}
Дано исходное 8-битное слово $D5_{16}$. Подпункты независимы. \textbf{По 3 балла за подпункт.}
\begin{enumerate}
    \item Логически сдвиньте вправо на 2; укажите биты, hex и беззнаковое значение.
    \item Арифметически сдвиньте вправо на 2; укажите биты, hex и знаковое значение.
    \item Циклически сдвиньте влево на 1; укажите биты и hex-код.
    \item Выполните знаковое расширение до 16 битов; укажите биты и hex-код.
    \item Прочитайте исходное слово как знаковое число, разделите на 4 с округлением к нулю и сравните с подпунктом 2.
\end{enumerate}

\section{В6. Кодирование float32 (15 баллов, 14 минут)}\label{task:v6}
\begin{enumerate}
    \item Представьте $-13.375_{10}$ в \texttt{float32}. Покажите нормализованную двоичную запись, знак, истинный и хранимый порядки, 23 бита дробной части. Запишите итоговые 32 бита по полям и восемь hex-цифр. \textbf{(9 баллов)}
    \item Декодируйте \texttt{0x40B80000} как \texttt{float32}. Выделите поля, найдите истинный порядок, значащую часть со скрытым битом и десятичное значение. \textbf{(6 баллов)}
\end{enumerate}

\section{В7. Классификация float32 (10 баллов, 10 минут)}\label{task:v7}
Для каждого слова определите знак, поле порядка и нулевое или ненулевое содержимое дробного поля, затем назовите класс значения:
\begin{enumerate}
    \item \texttt{0x80000000}; \textbf{(2 балла)}
    \item \texttt{0x7F800000}; \textbf{(2 балла)}
    \item \texttt{0x7FC00001}; \textbf{(2 балла)}
    \item \texttt{0x00400000}. Для этого слова дополнительно вычислите точное значение в виде степени двойки, показав формулу со значащей частью. \textbf{(4 балла)}
\end{enumerate}

\section{В8. Точность и округление (15 баллов, 12 минут)}\label{task:v8}
Все числа и промежуточные результаты хранятся в \texttt{float32}. Дано $h=2^{-24}$. Порядок операций соответствует скобкам; каждая операция округляется по правилу, указанному в начале варианта.
\begin{enumerate}
    \item Найдите машинный эпсилон как расстояние от 1 до следующего большего представимого числа. \textbf{(3 балла)}
    \item Вычислите $(1+h)+h$, показывая округление после каждого сложения. \textbf{(4 балла)}
    \item Вычислите $1+(h+h)$, показывая оба промежуточных шага. \textbf{(4 балла)}
    \item Объясните различие результатов через расстояния до соседних чисел и правило выбора чётной значащей части. Назовите проявившиеся свойства арифметики. \textbf{(4 балла)}
\end{enumerate}
Ответы допускаются в виде точных выражений со степенями двойки.

\section*{Последние 10 минут: проверка}
Проверьте разрядность слов и порядок флагов. Для \texttt{float32} проверьте длину полей, скрытый бит и округление промежуточных результатов. После завершения работы переходите к \hyperref[answers:mock13]{кратким ответам} и \hyperref[solutions:mock13]{подробным решениям с баллами}.
